\scanpage{75}
% archive foliation: 74
Les développements précédents étaient centrés
sur le groupe $H=GP(1,\mathbf{C})$ agissant sur $\mathbf{P}^{1}(\mathbf{C})$
(ou $GP(1)_{\mathbf{Z}}$ sur $\mathbf{P}^{1}_{\mathbf{Z}}\ldots$), et les invariants
différentiels, sous l'action de ce groupe, des
fonctions analytiques complexes (\uline{NB} \unsure{tout cela}
vaut aussi dans l'analytique réel, ou le
$C^{\infty}$, ou $C^3$).
On va passer rapidement en revue les développements
analogues pour d'autres groupes

a) \uline{Groupe des translations} $G_a$ \uline{agissant sur $\unsure{\underline{E}^{1}|\underline{A}^{1}}$}

On a l'opération (dans l'analytique complexe
p.ex.)
\[
d:\underline{\mathcal{O}}_X\longrightarrow\underline{\Omega}^{1}_{X/\mathbf{C}}
=\underline{\omega}_X
\]
qui \inserted{épi ($X$ lisse dim $1$) et}
s'annule exactement sur les fonctions
loc. constantes, donc vient
\[
A(f,g)=dg-df,\qquad
A:\underline{\mathcal{O}}_X\times\underline{\mathcal{O}}_X
\longrightarrow\underline{\omega}_X
\]
\margin{Ici $\underline{\mathcal{O}}_X$ \unsure{mauvais};
\unsure{prendre} $\mathcal{E}_X$ plutôt!}

on trouve l'analogue de a) b) dessus,
plus
\[
A(f,g+c)=A(f,g)\qquad c\in\mathbf{C}
\]
plus
\[
A(f,g)=A(f,h),\quad g(x)=h(x)\quad\Longrightarrow\quad
\underset{\text{au vois. de }x}{g=h}
\]
et
\[
g\longmapsto A(f,g)\qquad
\mathcal{E}_X\longrightarrow\underline{\omega}_X\quad\text{épi}
\]
(Cas $X$ lisse de dim $1$ \unsure{p. sûr|p. ass.}) ---

Ici pas la peine d'introduire de torseur,
\inserted{(sinon, il est trivial)}
\unsure{un} $=$ $A(f,g)$, \unsure{car} l'expression
des \unsure{lettres} \unsure{telles} qui en \unsure{occurrent|sortent}

